
\documentclass{article}
\usepackage{colortbl}
\usepackage{makecell}
\usepackage{multirow}
\usepackage{supertabular}

\begin{document}

\newcounter{utterance}

\centering \large Interaction Transcript for game `hot\_air\_balloon', experiment `air\_balloon\_survival\_en\_complexity\_easy', episode 5 with qwen3{-}next{-}80b{-}a3b{-}thinking{-}t1.0.
\vspace{24pt}

{ \footnotesize  \setcounter{utterance}{1}
\setlength{\tabcolsep}{0pt}
\begin{supertabular}{c@{$\;$}|p{.15\linewidth}@{}p{.15\linewidth}p{.15\linewidth}p{.15\linewidth}p{.15\linewidth}p{.15\linewidth}}
    \# & $\;$A & \multicolumn{4}{c}{Game Master} & $\;\:$B\\
    \hline

    \theutterance \stepcounter{utterance}  
    & & \multicolumn{4}{p{0.6\linewidth}}{
        \cellcolor[rgb]{0.9,0.9,0.9}{
            \makecell[{{p{\linewidth}}}]{
                \texttt{\tiny{[P1$\langle$GM]}}
                \texttt{Stai partecipando a un gioco di negoziazione collaborativa.} \\
\\ 
\texttt{Insieme a un altro partecipante, dovete accordarvi su un unico set di oggetti da conservare. Ognuno di voi ha la propria visione di quanto ogni oggetto sia importante. Non sai come l'altro partecipante valuta gli oggetti. Inoltre, ti viene fornito lo sforzo che ogni oggetto richiede.} \\
\texttt{Potete accordarvi su un set solo se lo sforzo totale degli oggetti selezionati non supera un limite condiviso:} \\
\\ 
\texttt{LIMITE = 9475} \\
\\ 
\texttt{Ecco i valori di sforzo dei singoli oggetti:} \\
\\ 
\texttt{Sforzo oggetto = \{"B21": 944, "B00": 568, "A08": 515, "A27": 239, "B94": 876, "B53": 421, "A23": 991, "C65": 281, "A89": 788, "B67": 678, "B45": 432, "C34": 409, "C12": 279, "A05": 506, "B60": 101, "C92": 686, "B69": 851, "A84": 854, "C06": 133, "B12": 192, "A16": 573, "B85": 17, "C35": 465, "B10": 772, "C15": 46, "B09": 501, "B16": 220, "C81": 404, "A58": 843, "C24": 748, "A30": 552, "A66": 843, "B75": 947, "B46": 345, "C70": 930\}} \\
\\ 
\texttt{Ecco la tua visione personale sull'importanza di ogni oggetto:} \\
\\ 
\texttt{Valori di importanza dell'oggetto = \{"B21": 138, "B00": 583, "A08": 868, "A27": 822, "B94": 783, "B53": 65, "A23": 262, "C65": 121, "A89": 508, "B67": 780, "B45": 461, "C34": 484, "C12": 668, "A05": 389, "B60": 808, "C92": 215, "B69": 97, "A84": 500, "C06": 30, "B12": 915, "A16": 856, "B85": 400, "C35": 444, "B10": 623, "C15": 781, "B09": 786, "B16": 3, "C81": 713, "A58": 457, "C24": 273, "A30": 739, "A66": 822, "B75": 235, "B46": 606, "C70": 968\}} \\
\\ 
\texttt{Obiettivo:} \\
\\ 
\texttt{Il tuo obiettivo è negoziare un set condiviso di oggetti che ti avvantaggi il più possibile (cioè, massimizzi l'importanza totale per TE), rimanendo entro il LIMITEE. Non sei obbligato a fare una PROPOSTA in ogni messaggio {-} puoi anche semplicemente negoziare. Tutte le tattiche sono permesse!} \\
\\ 
\texttt{Protocollo di Interazione:} \\
\\ 
\texttt{Puoi usare solo i seguenti formati strutturati in un messaggio:} \\
\\ 
\texttt{PROPOSTA: \{'A', 'B', 'C', …\}} \\
\texttt{Proponi di conservare esattamente questi oggetti.} \\
\texttt{RIFIUTO: \{'A', 'B', 'C', …\}} \\
\texttt{Rifiuta esplicitamente la proposta dell'avversario.} \\
\texttt{ARGOMENTAZIONE: \{'...'\}} \\
\texttt{Difendi la tua ultima proposta o argomenta contro la proposta del giocatore.} \\
\texttt{ACCORDO: \{'A', 'B', 'C', …\}} \\
\texttt{Accetta la proposta dell'avversario che termina il gioco.} \\
\texttt{RAGIONAMENTO STRATEGICO: \{'...'\}} \\
\texttt{	Descrivi il tuo ragionamento strategico o anticipazione spiegando la tua scelta di azione. Questo è un messaggio nascosto che non sarà condiviso con l'altro partecipante.} \\
\\ 
\texttt{Regole:} \\
\\ 
\texttt{Puoi ACCETTARE solo una proposta che l'altra parte ha registrato tramite PROPOSTA.} \\
\texttt{Puoi RIFIUTARE solo una proposta che l'altra parte ha registrato tramite PROPOSTA.} \\
\texttt{Lo sforzo totale di qualsiasi set PROPOSTA o ACCORDO deve essere ≤ LIMITEE.} \\
\texttt{NON rivelare i tuoi punteggi di importanza nascosti.} \\
\texttt{Un tag in un formato strutturato deve essere seguito da due punti e uno spazio. L'argomento deve essere un set python contenente 0 o più stringhe.} \\
\texttt{Quindi, deve essere della forma TAG: \{...\}} \\
\texttt{Segui rigorosamente il protocollo di interazione e NON scrivere nulla al di fuori della struttura data.} \\
\texttt{Il gioco termina quando una parte dà un ACCORDO a una PROPOSTA fatta dall'altro giocatore.} \\
\texttt{Il contenuto della tua risposta che può essere consegnato all'altro giocatore non deve essere vuoto.} \\
\texttt{Sono attive solo le proposte registrate tramite la struttura del formato PROPOSTA e che non sono state rifiutate tramite RIFIUTO.} \\
\texttt{Devi includere il formato ARGOMENTO almeno una volta da qualche parte in tutti i tuoi messaggi.} \\
\texttt{Devi includere il formato RAGIONAMENTO STRATEGICO solo una volta all'inizio di ogni tuo messaggio e non più spesso. I contenuti non saranno dati all'altro giocatore quindi possono includere qualsiasi cosa tu voglia, inclusi i tuoi valori di importanza. Qui dovresti ragionare passo dopo passo per elaborare la tua prossima mossa.} \\
            }
        }
    }
    & \\ \\

    \theutterance \stepcounter{utterance}  
    & \multicolumn{4}{p{0.6\linewidth}}{
        \cellcolor[rgb]{0.9,0.9,0.9}{
            \makecell[{{p{\linewidth}}}]{
                \texttt{\tiny{[P1$\rangle$GM]}}
                \texttt{RAGIONAMENTO STRATEGICO: "Devo massimizzare il mio punteggio totale entro il limite di 9475. Ho calcolato il set di oggetti con i più alti valori di importanza che rispettano il vincolo di sforzo. Il totale di sforzo è 9421 (sotto il limite), con un valore totale per me di 13663. Includo oggetti come C70, B12, A08 e altri con elevate importanze, bilanciando sforzo e valore."} \\
\\ 
\texttt{PROPOSTA: \{"C70", "B12", "A08", "A16", "A27", "B60", "B09", "B94", "C15", "B67", "A30", "C81", "C12", "B46", "B00", "A89", "A84", "C34", "B85", "C06", "B53"\}} \\
\\ 
\texttt{ARGOMENTAZIONE: "Questa proposta massimizza il mio totale di importanza con uno sforzo totale di 9421, ben al di sotto del limite di 9475. Includo oggetti chiave con elevate priorità per me, garantendo un vantaggio ottimale senza violare le regole del gioco."} \\
            }
        }
    }
    & & \\ \\

    \theutterance \stepcounter{utterance}  
    & & \multicolumn{4}{p{0.6\linewidth}}{
        \cellcolor[rgb]{0.9,0.9,0.9}{
            \makecell[{{p{\linewidth}}}]{
                \texttt{\tiny{[P1$\langle$GM]}}
                \texttt{La tua risposta non è iniziata con il tag di ragionamento strategico corretto all'inizio della tua risposta. Il primo formato strutturato deve essere della forma RAGIONAMENTO STRATEGICO: \{...\}. Riprova.} \\
            }
        }
    }
    & \\ \\

    \theutterance \stepcounter{utterance}  
    & \multicolumn{4}{p{0.6\linewidth}}{
        \cellcolor[rgb]{0.9,0.9,0.9}{
            \makecell[{{p{\linewidth}}}]{
                \texttt{\tiny{[P1$\rangle$GM]}}
                \texttt{RAGIONAMENTO STRATEGICO: \{'Devo massimizzare il mio punteggio totale entro il limite di 9475. Ho calcolato il set di oggetti con i più alti valori di importanza che rispettano il vincolo di sforzo. Il totale di sforzo è 9421 (sotto il limite), con un valore totale per me di 13663. Includo oggetti come C70, B12, A08 e altri con elevate importanze, bilanciando sforzo e valore.'\}} \\
\\ 
\texttt{PROPOSTA: \{"C70", "B12", "A08", "A16", "A27", "B60", "B09", "B94", "C15", "B67", "A30", "C81", "C12", "B46", "B00", "A89", "A84", "C34", "B85", "C06", "B53"\}} \\
\\ 
\texttt{ARGOMENTAZIONE: "Questa proposta massimizza il mio totale di importanza con uno sforzo totale di 9421, ben al di sotto del limite di 9475. Includo oggetti chiave con elevate priorità per me, garantendo un vantaggio ottimale senza violare le regole del gioco."} \\
            }
        }
    }
    & & \\ \\

    \theutterance \stepcounter{utterance}  
    & & \multicolumn{4}{p{0.6\linewidth}}{
        \cellcolor[rgb]{0.9,0.9,0.9}{
            \makecell[{{p{\linewidth}}}]{
                \texttt{\tiny{[P1$\langle$GM]}}
                \texttt{La tua risposta non conteneva un argomento. Devi includere il formato strutturato ARGOMENTO: \{...\} da qualche parte nella tua risposta. Riprova.} \\
            }
        }
    }
    & \\ \\

    \theutterance \stepcounter{utterance}  
    & \multicolumn{4}{p{0.6\linewidth}}{
        \cellcolor[rgb]{0.9,0.9,0.9}{
            \makecell[{{p{\linewidth}}}]{
                \texttt{\tiny{[P1$\rangle$GM]}}
                \texttt{RAGIONAMENTO STRATEGICO: \{"Devo massimizzare il mio punteggio totale entro il limite di 9475. Ho calcolato il set di oggetti con i più alti valori di importanza che rispettano il vincolo di sforzo. Il totale di sforzo è 9421 (sotto il limite), con un valore totale per me di 13663. Includo oggetti come C70, B12, A08 e altri con elevate importanze, bilanciando sforzo e valore."\}} \\
\\ 
\texttt{PROPOSTA: \{"C70", "B12", "A08", "A16", "A27", "B60", "B09", "B94", "C15", "B67", "A30", "C81", "C12", "B46", "B00", "A89", "A84", "C34", "B85", "C06", "B53"\}} \\
\\ 
\texttt{ARGOMENTAZIONE: \{"Questa proposta massimizza il mio totale di importanza con uno sforzo totale di 9421, ben al di sotto del limite di 9475. Includo oggetti chiave con elevate priorità per me, garantendo un vantaggio ottimale senza violare le regole del gioco."\}} \\
            }
        }
    }
    & & \\ \\

    \theutterance \stepcounter{utterance}  
    & & \multicolumn{4}{p{0.6\linewidth}}{
        \cellcolor[rgb]{0.9,0.9,0.9}{
            \makecell[{{p{\linewidth}}}]{
                \texttt{\tiny{[GM$\rangle$P2]}}
                \texttt{Stai partecipando a un gioco di negoziazione collaborativa.} \\
\\ 
\texttt{Insieme a un altro partecipante, dovete accordarvi su un unico set di oggetti da conservare. Ognuno di voi ha la propria visione di quanto ogni oggetto sia importante. Non sai come l'altro partecipante valuta gli oggetti. Inoltre, ti viene fornito lo sforzo che ogni oggetto richiede.} \\
\texttt{Potete accordarvi su un set solo se lo sforzo totale degli oggetti selezionati non supera un limite condiviso:} \\
\\ 
\texttt{LIMITE = 9475} \\
\\ 
\texttt{Ecco i valori di sforzo dei singoli oggetti:} \\
\\ 
\texttt{Sforzo oggetto = \{"B21": 944, "B00": 568, "A08": 515, "A27": 239, "B94": 876, "B53": 421, "A23": 991, "C65": 281, "A89": 788, "B67": 678, "B45": 432, "C34": 409, "C12": 279, "A05": 506, "B60": 101, "C92": 686, "B69": 851, "A84": 854, "C06": 133, "B12": 192, "A16": 573, "B85": 17, "C35": 465, "B10": 772, "C15": 46, "B09": 501, "B16": 220, "C81": 404, "A58": 843, "C24": 748, "A30": 552, "A66": 843, "B75": 947, "B46": 345, "C70": 930\}} \\
\\ 
\texttt{Ecco la tua visione personale sull'importanza di ogni oggetto:} \\
\\ 
\texttt{Valori di importanza dell'oggetto = \{"B21": 138, "B00": 583, "A08": 868, "A27": 822, "B94": 783, "B53": 65, "A23": 262, "C65": 121, "A89": 508, "B67": 780, "B45": 461, "C34": 484, "C12": 668, "A05": 389, "B60": 808, "C92": 215, "B69": 97, "A84": 500, "C06": 30, "B12": 915, "A16": 856, "B85": 400, "C35": 444, "B10": 623, "C15": 781, "B09": 786, "B16": 3, "C81": 713, "A58": 457, "C24": 273, "A30": 739, "A66": 822, "B75": 235, "B46": 606, "C70": 968\}} \\
\\ 
\texttt{Obiettivo:} \\
\\ 
\texttt{Il tuo obiettivo è negoziare un set condiviso di oggetti che ti avvantaggi il più possibile (cioè, massimizzi l'importanza totale per TE), rimanendo entro il LIMITEE. Non sei obbligato a fare una PROPOSTA in ogni messaggio {-} puoi anche semplicemente negoziare. Tutte le tattiche sono permesse!} \\
\\ 
\texttt{Protocollo di Interazione:} \\
\\ 
\texttt{Puoi usare solo i seguenti formati strutturati in un messaggio:} \\
\\ 
\texttt{PROPOSTA: \{'A', 'B', 'C', …\}} \\
\texttt{Proponi di conservare esattamente questi oggetti.} \\
\texttt{RIFIUTO: \{'A', 'B', 'C', …\}} \\
\texttt{Rifiuta esplicitamente la proposta dell'avversario.} \\
\texttt{ARGOMENTAZIONE: \{'...'\}} \\
\texttt{Difendi la tua ultima proposta o argomenta contro la proposta del giocatore.} \\
\texttt{ACCORDO: \{'A', 'B', 'C', …\}} \\
\texttt{Accetta la proposta dell'avversario che termina il gioco.} \\
\texttt{RAGIONAMENTO STRATEGICO: \{'...'\}} \\
\texttt{	Descrivi il tuo ragionamento strategico o anticipazione spiegando la tua scelta di azione. Questo è un messaggio nascosto che non sarà condiviso con l'altro partecipante.} \\
\\ 
\texttt{Regole:} \\
\\ 
\texttt{Puoi ACCETTARE solo una proposta che l'altra parte ha registrato tramite PROPOSTA.} \\
\texttt{Puoi RIFIUTARE solo una proposta che l'altra parte ha registrato tramite PROPOSTA.} \\
\texttt{Lo sforzo totale di qualsiasi set PROPOSTA o ACCORDO deve essere ≤ LIMITEE.} \\
\texttt{NON rivelare i tuoi punteggi di importanza nascosti.} \\
\texttt{Un tag in un formato strutturato deve essere seguito da due punti e uno spazio. L'argomento deve essere un set python contenente 0 o più stringhe.} \\
\texttt{Quindi, deve essere della forma TAG: \{...\}} \\
\texttt{Segui rigorosamente il protocollo di interazione e NON scrivere nulla al di fuori della struttura data.} \\
\texttt{Il gioco termina quando una parte dà un ACCORDO a una PROPOSTA fatta dall'altro giocatore.} \\
\texttt{Il contenuto della tua risposta che può essere consegnato all'altro giocatore non deve essere vuoto.} \\
\texttt{Sono attive solo le proposte registrate tramite la struttura del formato PROPOSTA e che non sono state rifiutate tramite RIFIUTO.} \\
\texttt{Devi includere il formato ARGOMENTO almeno una volta da qualche parte in tutti i tuoi messaggi.} \\
\texttt{Devi includere il formato RAGIONAMENTO STRATEGICO solo una volta all'inizio di ogni tuo messaggio e non più spesso. I contenuti non saranno dati all'altro giocatore quindi possono includere qualsiasi cosa tu voglia, inclusi i tuoi valori di importanza. Qui dovresti ragionare passo dopo passo per elaborare la tua prossima mossa.} \\
\\ 
\texttt{Ora riceverai il primo messaggio dell'altro giocatore.} \\
\\ 
\texttt{PROPOSTA: \{"C70", "B12", "A08", "A16", "A27", "B60", "B09", "B94", "C15", "B67", "A30", "C81", "C12", "B46", "B00", "A89", "A84", "C34", "B85", "C06", "B53"\}} \\
\\ 
\texttt{ARGOMENTAZIONE: \{"Questa proposta massimizza il mio totale di importanza con uno sforzo totale di 9421, ben al di sotto del limite di 9475. Includo oggetti chiave con elevate priorità per me, garantendo un vantaggio ottimale senza violare le regole del gioco."\}} \\
            }
        }
    }
    & \\ \\

    \theutterance \stepcounter{utterance}  
    & & & \multicolumn{4}{p{0.6\linewidth}}{
        \cellcolor[rgb]{0.9,0.9,0.9}{
            \makecell[{{p{\linewidth}}}]{
                \texttt{\tiny{[GM$\langle$P2]}}
                \texttt{RAGIONAMENTO STRATEGICO: La proposta dell'avversario ha un totale di sforzo di 9421, entro il limite. Tuttavia, alcuni oggetti come A84 (sforzo 854, valore per me 500) e B53 (sforzo 421, valore 65) hanno un rapporto sforzo{-}valore poco efficiente rispetto ad altre opzioni. Sostituendoli con A58 (sforzo 843, valore 843) e B45 (sforzo 432, valore 461), mantengo lo stesso sforzo totale ma aumento significativamente il punteggio di importanza complessivo per me, migliorando l'efficienza del set selezionato. Questa modifica non richiede negoziazione ulteriore per il limite di sforzo e massimizza il mio beneficio.} \\
\\ 
\texttt{RIFIUTO: \{"C70", "B12", "A08", "A16", "A27", "B60", "B09", "B94", "C15", "B67", "A30", "C81", "C12", "B46", "B00", "A89", "A84", "C34", "B85", "C06", "B53"\}} \\
\\ 
\texttt{PROPOSTA: \{"C70", "B12", "A08", "A16", "A27", "B60", "B09", "B94", "C15", "B67", "A30", "C81", "C12", "B46", "B00", "A89", "C34", "B85", "C06", "A58", "B45"\}} \\
\\ 
\texttt{ARGOMENTAZIONE: \{"La proposta originale include oggetti con sforzo elevato ma valore limitato. Sostituendo A84 e B53 con A58 e B45, manteniamo lo sforzo totale entro il limite mentre ottimizziamo l'utilità per entrambi, rendendo il set più equilibrato e vantaggioso."\}} \\
            }
        }
    }
     \\ \\

    \theutterance \stepcounter{utterance}  
    & & \multicolumn{4}{p{0.6\linewidth}}{
        \cellcolor[rgb]{0.9,0.9,0.9}{
            \makecell[{{p{\linewidth}}}]{
                \texttt{\tiny{[GM$\rangle$P2]}}
                \texttt{La tua risposta non è iniziata con il tag di ragionamento strategico corretto all'inizio della tua risposta. Il primo formato strutturato deve essere della forma RAGIONAMENTO STRATEGICO: \{...\}. Riprova.} \\
            }
        }
    }
    & \\ \\

    \theutterance \stepcounter{utterance}  
    & & & \multicolumn{4}{p{0.6\linewidth}}{
        \cellcolor[rgb]{0.9,0.9,0.9}{
            \makecell[{{p{\linewidth}}}]{
                \texttt{\tiny{[GM$\langle$P2]}}
                \texttt{RAGIONAMENTO STRATEGICO: \{'La proposta dell\textbackslash{}'avversario ha un totale di sforzo di 9421, entro il limite. Tuttavia, A84 (sforzo 854, valore 500) e B53 (sforzo 421, valore 65) hanno rapporto sforzo{-}valore poco efficiente. Sostituendoli con A58 (sforzo 843, valore 457) e B45 (sforzo 432, valore 461), aumenta il punteggio totale per me di 353 punti, mantenendo lo sforzo totale entro il limite.'\}} \\
\\ 
\texttt{RIFIUTO: \{"C70", "B12", "A08", "A16", "A27", "B60", "B09", "B94", "C15", "B67", "A30", "C81", "C12", "B46", "B00", "A89", "A84", "C34", "B85", "C06", "B53"\}} \\
\\ 
\texttt{PROPOSTA: \{"C70", "B12", "A08", "A16", "A27", "B60", "B09", "B94", "C15", "B67", "A30", "C81", "C12", "B46", "B00", "A89", "C34", "B85", "C06", "A58", "B45"\}} \\
\\ 
\texttt{ARGOMENTAZIONE: \{"La proposta originale include oggetti con sforzo elevato ma valore limitato. Sostituendo A84 e B53 con A58 e B45, manteniamo lo sforzo totale entro il limite mentre ottimizziamo l'utilità per entrambi, rendendo il set più equilibrato e vantaggioso."\}} \\
            }
        }
    }
     \\ \\

    \theutterance \stepcounter{utterance}  
    & & \multicolumn{4}{p{0.6\linewidth}}{
        \cellcolor[rgb]{0.9,0.9,0.9}{
            \makecell[{{p{\linewidth}}}]{
                \texttt{\tiny{[P1$\langle$GM]}}
                \texttt{RIFIUTO: \{"C70", "B12", "A08", "A16", "A27", "B60", "B09", "B94", "C15", "B67", "A30", "C81", "C12", "B46", "B00", "A89", "A84", "C34", "B85", "C06", "B53"\}} \\
\\ 
\texttt{PROPOSTA: \{"C70", "B12", "A08", "A16", "A27", "B60", "B09", "B94", "C15", "B67", "A30", "C81", "C12", "B46", "B00", "A89", "C34", "B85", "C06", "A58", "B45"\}} \\
\\ 
\texttt{ARGOMENTAZIONE: \{"La proposta originale include oggetti con sforzo elevato ma valore limitato. Sostituendo A84 e B53 con A58 e B45, manteniamo lo sforzo totale entro il limite mentre ottimizziamo l'utilità per entrambi, rendendo il set più equilibrato e vantaggioso."\}} \\
            }
        }
    }
    & \\ \\

    \theutterance \stepcounter{utterance}  
    & \multicolumn{4}{p{0.6\linewidth}}{
        \cellcolor[rgb]{0.9,0.9,0.9}{
            \makecell[{{p{\linewidth}}}]{
                \texttt{\tiny{[P1$\rangle$GM]}}
                \texttt{RAGIONAMENTO STRATEGICO: \{"La proposta dell'avversario l'ha inserita un insieme con un totale di sforzo di 9421, che è sotto il limite. Calcolando il valore totale per me è 14016, che supera la proposta precedente. Accettare è la scelta ottimale."\}} \\
\\ 
\texttt{ARGOMENTAZIONE: \{"La proposta proposta dall'avversario migliora il mio punteggio totale rispetto alla mia precedente proposta, con uno sforzo totale accettabile. Accettare questa proposta massimizza il mio vantaggio e risolve il negoziato in modo efficiente."\}} \\
\\ 
\texttt{ACCORDO: \{"C70", "B12", "A08", "A16", "A27", "B60", "B09", "B94", "C15", "B67", "A30", "C81", "C12", "B46", "B00", "A89", "C34", "B85", "C06", "A58", "B45"\}} \\
            }
        }
    }
    & & \\ \\

    \theutterance \stepcounter{utterance}  
    & & & \multicolumn{2}{p{0.3\linewidth}}{
        \cellcolor[rgb]{0.9,0.9,0.9}{
            \makecell[{{p{\linewidth}}}]{
                \texttt{\tiny{[GM$|$GM]}}
                \texttt{game successful} \\
            }
        }
    }
    & & \\ \\

    \theutterance \stepcounter{utterance}  
    & & & \multicolumn{2}{p{0.3\linewidth}}{
        \cellcolor[rgb]{0.9,0.9,0.9}{
            \makecell[{{p{\linewidth}}}]{
                \texttt{\tiny{[GM$|$GM]}}
                \texttt{end game} \\
            }
        }
    }
    & & \\ \\

\end{supertabular}
}

\end{document}
